\section{Chapter 3: Propositions and proofs}\label{sec:14-appendix-solutions:02-chapter-3:1}

\textbf{1. \texttt{theorem and\_\allowbreak{}comm\_\allowbreak{}ex \{P Q : Prop\} (h : P ∧ Q) : Q ∧ P}}

\begin{lstlisting}
theorem and_comm_ex {P Q : Prop} (h : P ∧ Q) : Q ∧ P :=
  ⟨h.right, h.left⟩
\end{lstlisting}

\texttt{h : P ∧ Q} has fields \texttt{h.left : P} and \texttt{h.right : Q}. A proof of \texttt{Q ∧ P} is just the pair with those two components in the opposite order.

\textbf{2. \texttt{theorem or\_\allowbreak{}comm\_\allowbreak{}ex \{P Q : Prop\} (h : P ∨ Q) : Q ∨ P}}

\begin{lstlisting}
theorem or_comm_ex {P Q : Prop} (h : P ∨ Q) : Q ∨ P :=
  match h with
  | Or.inl hp => Or.inr hp
  | Or.inr hq => Or.inl hq
\end{lstlisting}

\texttt{Or} has two constructors, so any proof \texttt{h : P ∨ Q} was built with one of them. \texttt{match} finds out which one, and the witness it carried. In the \texttt{Or.inl hp} case we have \texttt{hp : P}, and \texttt{Or.inr hp : Q ∨ P} uses it as the right disjunct. The \texttt{Or.inr hq} case works the same way, just mirrored.

\textbf{3. \texttt{theorem exists\_\allowbreak{}gt\_\allowbreak{}zero : ∃ n : Nat, n > 0}}

\begin{lstlisting}
theorem exists_gt_zero : ∃ n : Nat, n > 0 :=
  ⟨1, by decide⟩
\end{lstlisting}

An \texttt{∃}-proof is a witness (\texttt{1}) paired with a proof that it satisfies the predicate (\texttt{1 > 0}). Here \href{https://lean-lang.org/doc/reference/latest/Tactic-Proofs/Tactic-Reference/}{\texttt{decide}} handles that proof, since \texttt{1 > 0} on \texttt{Nat} is a decidable, closed proposition. We could also write \texttt{⟨1, Nat.one\_\allowbreak{}pos⟩} or \texttt{⟨1, rfl⟩} (since \texttt{1 > 0} unfolds to \texttt{0 < 1}, i.e. \texttt{Nat.succ 0 ≤ 1}, which is true by definition).
